\section{The fixed-word tree calculus}
\label{sec:foundations}

This section constructs the common space in which we compare ordinary
QAOA words and classical formulas.  It has three tasks: identify the
fixed-word SK value with a tree limit, express that limit as a bounded
quadratic form, and justify the controlled operations used in synthesis.
The same space will carry the finite Gaussian construction in
Section~\ref{sec:classical-optimization}.

We first isolate the transfer from an ordinary word to the tree.
A word always means a finite word with fixed real coefficients; none
of the statements below makes a uniform assertion in its length or
its coefficients. The initial vector $\vac$ below is the plus product
state on the finite light cone used to define the edge expectation.

\begin{proposition}[Fixed-word transfer]\label{prop:tree-transfer}
Let $W$ be an ordinary QAOA word, and apply the same word to the
$(D+1)$-regular tree with cost
$C_D=D^{-1/2}\sum_{\{i,j\}\in E}Z_iZ_j$, with each unordered edge
counted once.  Its expectation at an edge is
defined by restricting to the finite light cone of that edge.  Then
\begin{equation}\label{eq:tree-transfer}
 v_W=\lim_{D\to\infty}\frac{\sqrt D}{2}
       \bra{W_D\vac}Z_iZ_j\ket{W_D\vac}.
\end{equation}
Here $v_W$ is the disorder-averaged thermodynamic SK energy per vertex,
with the normalization in the introduction.  Both limits exist for
every fixed word and every fixed choice of its angles.
\end{proposition}

\begin{proof}
This is Theorem~1 of \cite{BFMVZ2022}, together with its
normalizations in Equations~(2.3), (2.6), and (6.1)--(6.2).
Its Equation~(3.11) gives the fixed-depth tree value, and
Equation~(6.3) identifies that value with the SK limit.
That paper uses the negative of our tree cost.
Conjugation by $X$ on one bipartition of the tree changes the
sign of every $Z_iZ_j$, preserves every mixer, and fixes the initial
vector.  It therefore changes the sign of the edge correlation as
well, giving precisely \eqref{eq:tree-transfer}.  The SK convention
already agrees with the positive Gaussian Hamiltonian in its
Equation~(6.1).  Consecutive gates of the same kind and inverses are
included by inserting zero-angle gates.
\end{proof}

We first construct the rooted space and identify its energy functional.
This also gives the Gaussian coordinates used by the classical
construction. We then prove the weighted convergence estimates that
realize ordinary words in this space and justify controlled gates inside
a fixed surrounding circuit.

\subsection{Pointed spaces and occupation embeddings}
\label{sec:pointed-spaces}

The number of children will tend to infinity while the number of
circuit layers remains fixed.  We therefore organize a symmetric
collection of child states by how many differ from the initial state.
The limiting space is a symmetric Fock space: its $m$-particle sector
records $m$ such child excitations, not $m$ physical spins.
The bosonic representation is related to the fixed-depth spin--boson
description of Boulebnane et al.~\cite{BKL2026}; the occupation,
insertion, and surrounding-circuit estimates used below are proved here.

Let $(S,\mu,\iota)$ be a probability space with a measure-preserving
involution.  The cases needed below are a one-point space and a
standard Gaussian, with $\iota(s)=-s$.  Extra independent marks may
be included in $S$.  Set
\[
 \mathsf k=\C^2\otimes L^2(S,\mu),\qquad
 \omega=\ket{+}\otimes\1.
\]
For any pointed space below, write $\langle T\rangle_0=
\bra{\vac}T\ket{\vac}$ for the vacuum mean and
$Q=1-|\vac\rangle\langle\vac|$ for its centered projection.
All seed variables act by multiplication.  In particular, a circuit
never changes their values.  The local symmetry is
$X\otimes(f\mapsto f\circ\iota)$.

For a pointed Hilbert space $(H,\vac)$ write
$H^\circ=H\ominus\C\vac$ and
$\Gamma_s(H^\circ)=\bigoplus_{m\geq0}(H^\circ)^{\otimes_s m}$.
At finite height define
\begin{equation}\label{eq:recursive-spaces}
 \cH_0=\mathsf k,\qquad
 \cH_{h+1}=\mathsf k\otimes\Gamma_s(\cH_h^\circ),\qquad
 \vac_{h+1}=\omega\otimes\Omega.
\end{equation}
The embedding from height $h$ to height $h+1$ appends empty
descendants: it is $\xi\mapsto\xi\otimes\Omega$ at height zero,
and thereafter is obtained by applying the symmetric Fock functor
to the preceding embedding.  These embeddings preserve the vacuum.
We identify the spaces with their images and use the completion of
their union when convenient.  Every construction in the proof takes
place at some finite height before this completion is used.

Write $R_\iota f=f\circ\iota$.  The joint spin/seed parity on these
spaces is the vacuum-preserving involution
\[
 \mathscr P_0=X\otimes R_\iota,\qquad
 \mathscr P_{h+1}=(X\otimes R_\iota)\otimes
       \Gamma_s(\mathscr P_h|_{\cH_h^\circ}).
\]
These involutions respect the height embeddings and fix the vacuum,
so every odd observable has zero vacuum mean.  We suppress the height
in $\mathscr P$ when the space is understood.

The analogous invariant space on a rooted $D$-ary tree is
\[
 \cH_{D,0}=\mathsf k,\qquad
 \cH_{D,h+1}=\mathsf k\otimes\operatorname{Sym}^D(\cH_{D,h}).
\]
The symmetry here permutes children together with their entire
subtrees and seed coordinates.  It does not discard the seeds.  For
any pointed space $H$, the orthogonal decomposition by the number
of nonvacuum tensor factors gives an isometry
\begin{equation}\label{eq:occupation-embedding}
 \mathcal I_D:\operatorname{Sym}^D(H)\longrightarrow
       \bigoplus_{m=0}^D(H^\circ)^{\otimes_s m}.
\end{equation}
On a symmetric $m$-particle tensor, the inverse distributes its $m$
nonvacuum factors over the $D$ slots with coefficient
$\binom Dm^{-1/2}$.  Applying this isometry recursively gives
isometries $\mathcal I_{D,h}:\cH_{D,h}\to\cH_h$.  Their range
projections tend strongly to the identity at each fixed height.
Indeed, a finite-particle vector involving finitely many child
vectors is approximated inductively by vectors in these ranges.
The finite-tree parity is the tensor product of the local
involutions over its vertices; the occupation embeddings
intertwine it with $\mathscr P_h$.

Creation is linear in its argument, annihilation is antilinear, and
\[
 q(d)=a^*(d)+a(d),\qquad
 [q(d),q(e)]=2\ii\im\braket{d}{e}.
\]
Conjugation by the recursive parity sends the child field $q(d)$
to $q(\mathscr P_h d)$; an odd child direction therefore gives
an odd field at the next height.
The fields are first defined on finite-particle vectors and then
closed.  Each $q(d)$ is essentially self-adjoint there, as follows
either from the one-mode Schr\"odinger representation or from the
convergence of its exponential series on finite-particle vectors.
We use $\mathcal W(d)=e^{\ii q(d)}$.  In particular,
\begin{equation}\label{eq:weyl-shift}
 \mathcal W(d)^*q(v)\mathcal W(d)
       =q(v)-2\im\braket{v}{d}.
\end{equation}

\subsection{Root insertions and edge energy}
\label{sec:root-energy}

At finite degree, a rooted vector records the response of a root
observable to a circuit. Let $U$ be a circuit on a finite subtree and
$U_j$ its restriction to child subtree $j$. Factoring the child
operations defines the root interaction-picture unitary $u$ by
\[
 U=\left(\bigotimes_{j\text{ child}}U_j\right)u.
\]
For a root Pauli observable $P$, its insertion vector is
\[
 f_P=U^*PU\vac=u^*Pu\vac,
\]
because the child operations commute with $P$. For a word evaluated
on a neighborhood containing both endpoint light cones, the pairing
of its two $Z$ insertions gives the expectation of $Z_iZ_j$.
We will prove convergence of ordinary insertion vectors
after establishing the weighted estimates below. First we identify
the energy of any convergent family of rooted vectors with bounded
expected occupation; this geometric step needs no circuit estimates.

The insertion vector lives at one root, whereas an edge joins two
rooted neighborhoods.  The following operator moves one centered
vector down to a single child while placing the new root in its
initial state.  It will turn the edge pairing into a quadratic form
on one common space.

Define the centered root shift by
\begin{equation}\label{eq:root-shift}
 \cL:\cH^\circ\longrightarrow\cH^\circ,
 \qquad \cL v=\omega\otimes a^*(v)\Omega.
\end{equation}
At height $h$ its range lies at height $h+1$.  The definitions
are compatible with the height embeddings, so $\cL$ is an
isometry on the centered direct limit.  Extend it by
$\cL\vac=0$ when writing $\cL^*$; this extension is a partial
isometry on the whole pointed space.  In particular $\cL^*$
extracts the root-vacuum one-particle component, with a constant
root seed.

\begin{proposition}[Energy as a root shift]\label{prop:energy}
Let $f_D$ be centered, root-invariant, finite-radius vectors on
the $(D+1)$-regular tree.  Root invariance means that the vector
itself is fixed by every rooted permutation of complete child
subtrees, including their seed coordinates, at every vertex of
its finite neighborhood.  Suppose their occupation embeddings
converge strongly to $v$ at a fixed finite height and their
expected top-generation occupations are bounded uniformly in
$D$.  Place the same rooted vector at the endpoints of an edge,
calling the resulting vectors $f_{D,i}$ and $f_{D,j}$.  Then
\begin{equation}\label{eq:energy-shift}
 \lim_{D\to\infty}\frac{\sqrt D}{2}
       \braket{f_{D,i}}{f_{D,j}}
       =\cE(v):=\re\braket{v}{\cL v}.
\end{equation}
For any centered $v,w$,
\begin{equation}\label{eq:energy-continuity}
 |\cE(v)-\cE(w)|
       \leq(\|v\|+\|w\|)\|v-w\|.
\end{equation}
\end{proposition}

\begin{proof}
We give the exact finite-degree decomposition.  Put $M=D+1$,
let $E$ be the pointed child space, and write $N_M$ for the
number of nonvacuum children in
$\mathsf k\otimes\operatorname{Sym}^M(E)$.  Let $R_M$ contract
a distinguished child against its vacuum, and let
\[
 J_M\xi=\omega\otimes\frac1{\sqrt M}
 \sum_{j=1}^M\vac^{\otimes(j-1)}\otimes\xi
                         \otimes\vac^{\otimes(M-j)}.
\]
The latter is an isometry from $E^\circ$.  Under occupation
embeddings, $J_M^*=\cL^*$ and
\begin{equation}\label{eq:vacuum-deletion}
 R_M=\sqrt{1-N/M},\qquad
 \|R_Mf-f\|^2\leq M^{-1}\bra fN_M\ket f.
\end{equation}
Indeed, in the $m$-particle sector the fraction of placements
surviving deletion is
$\binom{M-1}m/\binom Mm=1-m/M$.

Delete the distinguished edge and let $P_A,P_B$ denote the
vacuum projections on its two half-trees.  Decompose $f_{D,i}$
with $P_B+(1-P_B)$ and $f_{D,j}$ with $P_A+(1-P_A)$.
The term with both vacuum projections is zero: the surviving
half-tree vectors are centered.  If just one side is excited,
the other side being vacuum forces the root and all the other
branches of the first vector to be vacuum.  With
$s_D=J_M^*f_D$ and $r_D=R_Mf_D$, the two such terms are
$\braket{s_D}{r_D}/\sqrt M$ and its conjugate.  Heights in
this pairing are matched by adding empty descendants.  Each
remaining excited-side vector has squared norm
$\bra{f_D}N_M\ket{f_D}/M$, by symmetry among branches.  Thus
\begin{equation}\label{eq:exact-edge-decomposition}
 \braket{f_{D,i}}{f_{D,j}}
 =\frac2{\sqrt M}\re\braket{s_D}{r_D}+\varepsilon_D,
 \qquad
 |\varepsilon_D|\leq\frac{\bra{f_D}N_M\ket{f_D}}M.
\end{equation}
After multiplication by $\sqrt D/2$ the error tends to zero.
Equation~\eqref{eq:vacuum-deletion} and strong convergence show
that the first term tends to
$\re\braket{\cL^*v}{v}=\re\braket{v}{\cL v}$.
This proves \eqref{eq:energy-shift}.

Finally, \eqref{eq:energy-continuity} follows by adding and subtracting
$\braket{w}{\cL v}$ and using $\|\cL\|=1$.
\end{proof}

The occupation hypothesis in Proposition~\ref{prop:energy}
is a sufficient hypothesis for the displayed proof; strong
convergence alone does not bound the remainder in
\eqref{eq:exact-edge-decomposition}.  Lemma~\ref{lem:ordinary-updates} and
Theorem~\ref{thm:gate-calculus} will supply that hypothesis for
the quantum gates used here.

\begin{remark}[Classical rotations]\label{rem:classical-rotation-energy}
The classical sector consists of vectors whose spin is $\ket{+}$
at every vertex and whose coefficients depend only on seeds.
If a $Y$ rotation whose angle depends only on seeds produces an insertion
$v=\ket{+}m+\ket{-}n$ with seed-only coefficients in all
descendants, then $\cE(v)=\cE(\ket{+}m)$.  Indeed every
$\cL v$ has root spin $\ket{+}$, while
$\cL(\ket{-}n)$ has a child spin excitation and is orthogonal
to the classical sector.  Seed-only here permits dependence
on descendant seeds; it does not mean dependence only on the
root seed.
\end{remark}

\subsection{Classical Gaussian positions and descendant formulas}
\label{subsec:classical-gaussian-sector}

The root-shift functional also applies to classical vectors, before
they are compiled into quantum gates.  We now give a common Gaussian
realization for both classical constructions in this paper.  This is
a definition on the limiting pointed space; a finite-degree energy
limit still requires the hypotheses of Proposition~\ref{prop:energy}.

Keep every quantum spin in $\ket{+}$ and use real coefficients in the
seed and symmetric Fock factors.  This gives the real classical
subspace.  For a real centered classical child vector $f$, let
$\mathsf G(f)=q(f)$ be the position operator on the root's child
Fock factor.  For any finite list $f_1,\ldots,f_k$, these positions
strongly commute and have joint centered Gaussian vacuum law with
covariance $\braket{f_i}{f_j}$.  Indeed, choose a real orthonormal
basis of their finite-dimensional span.  The position operators of
these orthogonal modes act on separate Fock factors, and
\[
 \left\langle
       e^{\ii\sum_{j=1}^k t_j\mathsf G(f_j)}
 \right\rangle_0
 =\exp\left(-\frac12\left\|\sum_{j=1}^kt_jf_j\right\|^2\right).
\]
In particular, orthonormal directions give independent standard
normal positions: the variance convention is $\|f\|^2$, with no
additional factor of two.  The positions act trivially on the local
seed factor.  Their joint vacuum law is therefore independent of the
receiving root's seed, even when a child direction itself depends on
the seed at its own root.

We identify a scalar function of these commuting variables with the
vector obtained by applying that function to the vacuum.  Under this
identification,
\begin{equation}
 \norm{\bigl(\mathsf G(f)-\mathsf G(g)\bigr)\vac}
      =\norm{f-g}.
 \label{eq:classical-gaussian-isometry}
\end{equation}
All fields are realized on the same pointed space, so this equality
also specifies their coupling.  The empty-descendant embeddings in
\eqref{eq:recursive-spaces} preserve inner products, commute with
passing to the corresponding Fock fields, and keep the local seed
factor fixed.  Thus directions constructed at different finite
heights can first be embedded in one common height without changing
their joint Gaussian law.

\begin{definition}[Finite descendant formula]
\label{def:finite-descendant-formula}
A finite descendant formula is a finite expression built recursively
from a local seed and finitely many Gaussian positions of centered
formulas one level below, using pointwise scalar functions.  The
functions are required to be measurable and to produce square-integrable
vectors whenever the formula is used as a vector.  Additional empty
levels identify formulas of different radii.  Its finite-degree version
replaces every Gaussian position by the normalized sum of the
corresponding independent child variables.
\end{definition}

The definition by itself is a limiting-space construction, not a
finite-degree convergence theorem.  Later approximation arguments
impose boundedness or polynomial-growth conditions on the pointwise
maps as needed.  Because $\cL$ places a constant seed at the new root,
$\cL^*$ extracts precisely the first Gaussian chaos with constant
receiving-root seed.  If $e_1,\ldots,e_k$ are orthonormal real centered
child vectors and $F$ is a square-integrable function of
$\mathsf G(e_1),\ldots,\mathsf G(e_k)$ alone, this gives
\begin{equation}\label{eq:classical-first-chaos}
 \cL^*F
   =\sum_{j=1}^k
       \E\bigl[F\,\mathsf G(e_j)\bigr]e_j.
\end{equation}
Here $F$ may have a nonzero constant component, which $\cL^*$
annihilates.  The identity follows by the orthogonal projection onto
the first chaos in these Gaussian coordinates; the remaining
orthogonal child modes and all higher chaoses have zero projection.
This is the normalization used in the primary classical construction.

The rooted space, Gaussian positions, and energy functional now give
the setting for the classical construction in
Section~\ref{sec:classical-optimization}. The rest of this section
establishes convergence and continuity of the quantum operations that
will implement it.

\subsection{Weights and the topology of convergence}
\label{sec:smooth-topology}

The fields just defined are unbounded.  To pass their products and
parameter derivatives through limits, we need a common domain and
bounds that control both seed moments and excitations at every level.
These recursive weights serve that purpose; they are distinct from
the top-generation number weight used in the seed-removal argument.
If all local-seed coefficients under consideration are uniformly
bounded, take $W_0=1$ on $\mathsf k$.  For coefficients of polynomial
growth in a Gaussian seed, start with multiplication by $1+|s|$
and make it block diagonal at $\omega$:
\[
 W_0=P_0+Q_0(1+|s|)Q_0,
 \qquad P_0=\ket{\omega}\bra{\omega},\quad Q_0=1-P_0.
\]
Compression is understood on its natural closed domain.  This is
a positive self-adjoint operator bounded below by one.  It differs
from multiplication by $1+|s|$ by a bounded finite-rank operator
whose range consists of vectors with all seed moments.  Consequently
their positive-integer power domains coincide, with equivalent
graph norms.  For completeness, this last assertion follows by
induction from $A^m-B^m=\sum_{j=0}^{m-1}A^j(A-B)B^{m-1-j}$ on the
common core; every occurrence of $A-B$ has smooth finite-dimensional
range, and the resulting estimates close in both directions.

Since $W_0\omega=\omega$, we can recursively set
\begin{equation}\label{eq:recursive-weights}
 \Lambda_h=W_h|_{\cH_h^\circ},\qquad
 W_{h+1}=W_0\otimes1+1\otimes\dd\Gamma(\Lambda_h).
\end{equation}
Thus $W_h\vac_h=\vac_h$ and $W_h\geq1$.  On a finite tree use
$W_{D,h+1}=W_0+\sum_{j=1}^D Q_jW_{D,h}^{(j)}Q_j$.
The occupation embeddings intertwine these weights exactly.  This
also shows that their range projections commute with every power
of $W_h$.

Write
\[
 \mathscr S_h=\bigcap_{m\geq0}\operatorname{Dom}(W_h^m),
 \qquad \|\xi\|_m=\|W_h^m\xi\|.
\]
Finite-particle tensors of smooth child vectors and local vectors
with all seed moments are dense in every one of these graph norms.
One proof first cuts off the spectral variables of $W_0$ and each
$\Lambda_h$, then cuts off particle number, and finally approximates
in the resulting bounded spectral subspaces by algebraic tensors.
The finite-degree projections give simultaneous approximants in all
fixed graph norms.

When comparing unitary maps on the varying embedded spaces,
extend each embedded unitary by the identity on the orthogonal
complement of its range projection.  These are genuine
unitaries on the common Hilbert space.  The projections commute
with the weights and converge strongly to one in every fixed
graph norm on smooth vectors, so this convention preserves
all the bounds and limits.  Parameter derivatives of the
extensions vanish on that complement.

We will repeatedly use the following precise convention.  An
operator family $F_D(x)$ has \emph{uniform smooth bounds} if, for
every nonnegative integer $m$ and multi-index $\alpha$, there are
$r,k,C$, independent of $D$ and $x$, such that
\begin{equation}\label{eq:smooth-family-bound}
 \|W_D^m\partial_x^\alpha F_D(x)\xi\|
 \leq C(1+|x|)^k\|W_D^r\xi\|,
 \qquad \xi\in\mathscr S_D.
\end{equation}
The same condition is required of the adjoint family.  Constants
may depend on the fixed expression, its finite list of channels,
the height, $m$, and $\alpha$.  Dependence on a further limiting
parameter is explicitly stated when relevant.  Smooth convergence
of operator families means strong convergence on smooth inputs
of both the families and their adjoints, including all specified
parameter jets, with these uniform bounds.  The limits of the
adjoint jets must be the adjoints of the corresponding limiting
jets on the common domain; bounded adjoint jets alone do not
imply their convergence.  For vector families only the vector
convergence is required.  This definition also permits triangular
arrays of one-copy spaces under the embeddings above.

Two elementary consequences will be useful.  Products of a fixed
number of such families have the same bounds, with a possibly
larger loss of input powers.  Also, strong convergence together
with bounded higher graph norms gives convergence in every fixed
lower graph norm.  The latter follows from the spectral inequality
\begin{equation}\label{eq:weight-interpolation}
 \|W^m\xi\|\leq R^m\|\xi\|
       +R^{m-M}\|W^M\xi\|\qquad(M>m, R\geq1),
\end{equation}
by taking the strong limit first and $R\to\infty$ second.
Consequently it is enough to prove convergence on the dense domain
just described, and then use the uniform estimates.  This is the
meaning of convergence between different finite-degree spaces.

\begin{lemma}[Weyl estimates]\label{lem:weyl-bounds}
Let $N$ be the particle number on $\Gamma_s(H^\circ)$.  For every
integer $m\geq0$ there is $C_m$ such that
\begin{align}
 \|(1+N)^m\mathcal W(d)\xi\|
 &\leq C_m(1+\|d\|)^{2m}\|(1+N)^m\xi\|,
       \label{eq:weyl-number-bound}\\
 \|(\mathcal W(d)-1)\xi\|
 &\leq2\|d\|\|(N+1)^{1/2}\xi\|.
       \label{eq:weyl-small-bound}
\end{align}
For the weight $W=W_0+\dd\Gamma(\Lambda)$, the same conclusion
holds with a polynomial in finitely many $\|\Lambda^j d\|$ in
place of $(1+\|d\|)^{2m}$ and with a finite loss of input powers.
All statements hold for spin-controlled displacements as well.
\end{lemma}

\begin{proof}
On finite-particle vectors,
$\|a(d)\xi\|\leq\|d\|\|N^{1/2}\xi\|$ and
$\|a^*(d)\xi\|\leq\|d\|\|(N+1)^{1/2}\xi\|$.
The fundamental theorem of calculus applied to
$e^{\ii t q(d)}\xi$ proves \eqref{eq:weyl-small-bound}.
Conjugation gives
\[
 \mathcal W(d)^*N\mathcal W(d)=N+q(\ii d)+\|d\|^2.
\]
Commuting $N$ through a field changes that field to another field
with the same norm of its argument.  Expanding the $m$th power
and applying the creation and annihilation bounds therefore proves
\eqref{eq:weyl-number-bound}.  For the weighted version use
\[
 \mathcal W(d)^*\dd\Gamma(\Lambda)\mathcal W(d)
 =\dd\Gamma(\Lambda)+q(\ii\Lambda d)
       +\braket{d}{\Lambda d},
\]
and $[\dd\Gamma(\Lambda),a^*(d)]=a^*(\Lambda d)$.
Every expansion uses only finitely many weighted norms of $d$.
Approximation in graph norms proves the identities and estimates
on the stated domains.  A spin-controlled displacement is the
sum of two ordinary displacements on orthogonal spin subspaces.
\end{proof}

\subsection{Normalized sums and a unitary central limit}
\label{sec:unitary-clt}

A normalized sum creates or removes one child excitation in the
limit.  A product of small unitaries has the corresponding field
exponential as its limit, but may also retain a scalar second-order
term.  We prove both statements with the weighted bounds needed for
later substitution into a fixed circuit.

\begin{lemma}[Sums]\label{lem:centered-sums}
Let $T_D$ and $T_D^*$ have uniform smooth bounds on one-copy
spaces, and suppose their smooth limits exist.  If
$\braket{\vac}{T_D\vac}=0$, then under the occupation embeddings
\begin{equation}\label{eq:centered-sum-limit}
 D^{-1/2}\sum_{j=1}^D T_D^{(j)}
 \longrightarrow a^*(T\vac)+a(T^*\vac)
\end{equation}
smoothly.  Without centering,
$D^{-1}\sum_jT_D^{(j)}\to\braket{\vac}{T\vac}\,1$.
Both families have uniform smooth bounds.
\end{lemma}

\begin{proof}
Put $g=QT_D\vac$, $h=QT_D^*\vac$, and $A=QT_DQ$.
For centered $T_D$ its normalized sum has the exact occupation
representation
\begin{equation}\label{eq:sum-decomposition}
 a^*(g)\sqrt{1-N/D}
 +\sqrt{1-N/D}\,a(h)+D^{-1/2}\dd\Gamma(A),
\end{equation}
where factors beyond the truncation are zero.  On the $m$-particle
sector the two square roots express, respectively, the numbers
$D-m$ of vacant and $D-m+1$ of potentially filled slots.
The formula first holds on algebraic tensors and then on smooth
vectors.  Creation and annihilation have uniform weighted bounds
by the proof of Lemma~\ref{lem:weyl-bounds}.  For the remaining
term, a one-copy bound between powers of $\Lambda$ implies a bound
for its sum between powers of $1+\sum\Lambda_j$: apply the
one-copy bound to each occupied slot, use
$\Lambda_j\leq\sum\Lambda_j$, and pay at most one additional
power for the number of occupied slots.  This argument works at
any desired output power.  It proves the required bounds for
\eqref{eq:sum-decomposition}.  On a fixed-particle smooth tensor
the last term tends to zero and the square roots tend to one;
density and \eqref{eq:weight-interpolation} finish the proof.
For an average the same decomposition has off-diagonal factor
$D^{-1/2}$ and complementary-block factor $D^{-1}$, while its
vacuum block tends to the stated scalar.
\end{proof}

The sum lemma determines first-order fields.  Controlled gates also
require products of $D$ factors whose arguments are of size
$D^{-1/2}$.  The next lemma keeps their second-order vacuum contribution;
it cannot in general be recovered from the first-order field alone.

\begin{lemma}[Products of small unitaries]\label{lem:unitary-clt}
Let $K_D(x)$, $x\in\R^b$, be unitary, satisfy $K_D(0)=1$, and
have uniform smooth bounds for all jets and adjoint jets.  Suppose
these jets converge smoothly and
$\braket{\vac}{\partial_\alpha K_D(0)\vac}=0$ for every first
derivative.  Write
\[
 K_\alpha=\partial_\alpha K(0)=\ii T_\alpha,
 \qquad K_{\alpha\lambda}=\partial_\alpha\partial_\lambda K(0).
\]
The $T_\alpha$ are symmetric on the smooth domain.  For every
fixed $a\in\R^b$,
\begin{equation}\label{eq:unitary-product-limit}
 K_D(a/\sqrt D)^{\otimes D}\longrightarrow
 \exp\left\{\ii\sum_\alpha a_\alpha q(T_\alpha\vac)
 +\frac12\sum_{\alpha,\lambda}a_\alpha a_\lambda
 \left\langle K_{\alpha\lambda}
           +\tfrac12\{T_\alpha,T_\lambda\}\right\rangle_0
 \right\}.
\end{equation}
The scalar in this exponent is imaginary.  The convergence holds
smoothly with all $a$-derivatives, locally uniformly in $a$, and
the product has polynomial smooth bounds in $a$, uniform in $D$.
The assertion also holds with an additional finite list of
unscaled parameters $\theta$: use
$K_D(a/\sqrt D;\theta)^{\otimes D}$, assume
$K_D(0;\theta)=1$ and first-derivative centering for every
$\theta$, and assume the hypotheses jointly for all parameter
jets, uniformly on compact $\theta$-sets.  The conclusion then
includes all $\theta$-jets and the same compact uniformity.
\end{lemma}

\begin{proof}
We first identify the candidate limit from exponential vectors;
then we prove the bounds needed for smooth convergence.  If $x,z\perp\vac$
are smooth, the vectors
$(\vac+x/\sqrt D)^{\otimes D}$ converge to
$\varepsilon(x)=\bigoplus_{m\geq0}x^{\otimes m}/\sqrt{m!}$,
and similarly for $z$.  The coefficient of $D^{-1}$ in the
one-copy matrix element is
\begin{align*}
 c={}&\braket{z}{x}
 +\sum_\alpha a_\alpha
       \bigl(\braket{z}{K_\alpha\vac}
                      +\braket{\vac}{K_\alpha x}\bigr)
 +\frac12\sum_{\alpha,\lambda}a_\alpha a_\lambda
                        \langle K_{\alpha\lambda}\rangle_0.
\end{align*}
Taylor's formula makes the remaining error $o(D^{-1})$, also
for converging triangular arrays of smooth $x,z$.  The full
pairing consequently tends to $e^c$.  For
$t=\sum_\alpha a_\alpha T_\alpha\vac$, the exponential-vector
pairing of $e^{\ii q(t)}$ has an additional term $-\|t\|^2/2$.
Since
$\|t\|^2=\sum_{\alpha,\lambda}a_\alpha a_\lambda
\langle\{T_\alpha,T_\lambda\}/2\rangle_0$,
the correction in \eqref{eq:unitary-product-limit} is exactly
the required one.  Twice differentiating $K(x)^*K(x)=1$ at zero
shows that its real part is zero. The displayed limit is thus unitary.

We next prove uniform bounds for the product
$U_D(a)=K_D(a/\sqrt D)^{\otimes D}$. For one child let $A=Q\Lambda Q$ and
$B_D(a)=K_D(a/\sqrt D)^*A K_D(a/\sqrt D)$.  Taylor's formula
at zero, with the assumed bounds, gives
\begin{align*}
 |\langle B_D(a)\rangle_0|&\leq D^{-1}P(a),\\
 \|\Lambda^m QB_D(a)\vac\|
  +\|\Lambda^m QB_D(a)^*\vac\|&\leq D^{-1/2}P_m(a),
\end{align*}
and the complementary block has polynomial smooth bounds.  Here
and below $P_m$ denotes some fixed nonnegative polynomial in
$1+|a|$, whose degree is allowed to depend on $m$.
The same vacuum-row decomposition as in
\eqref{eq:sum-decomposition}, now applied to the unnormalized
sum of the $B_D(a)$, shows that
\[
 U_D(a)^*WU_D(a)=W_0+\sum_jB_D(a)^{(j)}
\]
has polynomial bounds from any sufficiently high input weight
to any prescribed output weight, uniformly in $D$.
Composing this estimate a fixed number of times bounds its powers.
Since $W^mU_D=U_D(U_D^*WU_D)^m$, this proves the assertion for
$U_D$ itself; the adjoint is treated identically.  All identities
may first be made on smooth tensors, which the product preserves,
and then closed using these very estimates.

For derivatives,
\[
 U_D(a)^*\partial_\alpha U_D(a)
 =D^{-1/2}\sum_j
 \bigl(K_D^*\partial_\alpha K_D\bigr)(a/\sqrt D)^{(j)}.
\]
The one-copy vacuum mean is $O(D^{-1/2}P(a))$, since it vanishes
at zero.  Lemma~\ref{lem:centered-sums}, with this scalar part
retained, gives a uniform bound.  Further differentiation produces
products of finitely many such sums.  A sum carrying $k\geq2$
derivative-scaling factors is bounded by the average estimate
times $D^{1-k/2}\leq1$.  Induction proves all jet bounds.

Weak convergence on the total set of exponential
vectors, together with equality of the limiting norms, proves
strong convergence.  Higher-weight convergence follows from
\eqref{eq:weight-interpolation}; derivative convergence follows
from finite differences and the uniform bounds on one further
derivative.  This also gives local uniformity in $a$.
Since the strong limit is unitary, the adjoints converge
strongly as well: for unitaries $U_D,U$, one has
$\|(U_D^*-U^*)\xi\|=\|\xi-U_DU^*\xi\|$.
The adjoint jet bounds already proved, followed by the same
finite-difference argument, give convergence of adjoint jets.

To verify the unscaled-parameter assertion, the only additional
estimate concerns a logarithmic derivative
$A_\nu(x;\theta)=K(x;\theta)^*\partial_{\theta_\nu}K(x;\theta)$.
It satisfies $A_\nu(0;\theta)=0$.  Moreover
\[
 \left\langle\partial_{x_\alpha}A_\nu(0;\theta)\right\rangle_0
 =\partial_{\theta_\nu}
       \left\langle\partial_{x_\alpha}K(0;\theta)\right\rangle_0
 =0.
\]
Its vacuum entry at $a/\sqrt D$ is therefore
$O(D^{-1}P(a))$, its vacuum row and column are
$O(D^{-1/2}P(a))$, and its complementary block is
$O(D^{-1/2}P(a))$ with smooth bounds.  The unnormalized
sum of its $D$ copies consequently has uniform smooth
bounds by the decomposition used above.  Every further
$\theta$-derivative has the same two vanishings, and mixed
$a$-derivatives are treated by the normalized-sum or average
estimates.  This proves all joint jet bounds.  Pointwise
convergence for fixed $\theta$, finite differences, and the
adjoint argument just given prove the parametric conclusion.
\end{proof}

\begin{remark}\label{rem:unitary-clt-domains}
The anticommutator in \eqref{eq:unitary-product-limit} is a
quadratic form on smooth vectors; its vacuum expectation is
unambiguous because all jets preserve that domain.  Essential
self-adjointness of the one-copy derivative $T_\alpha$ is not
needed: only the vector $T_\alpha\vac$ enters the self-adjoint
Fock field.  This avoids an unnecessary domain hypothesis.
\end{remark}

\subsection{Smooth limits for ordinary root maps}
\label{sec:ordinary-root-calculus}

We now prove that the ordinary root maps introduced in
Section~\ref{sec:root-energy} have smooth limits, and hence that their
limiting insertions have the root-shift energy computed there.
At height $h$, write $u_h$ for the limiting root map and put
$v_{P,h}=u_h^*Pu_h\vac_h$ for a root Pauli observable $P$.
The following induction constructs these maps from the identity.
In the limiting cost update below, the field direction at height $h$ is the
preceding word's insertion on a child subtree at height $h-1$:
\[
 v_{Z,h-1}=u_{h-1}^*Zu_{h-1}\vac_{h-1},\qquad
 u_h\longmapsto e^{-\ii tZq(v_{Z,h-1})}u_h.
\]
Here $q(v_{Z,h-1})$ acts on the child Fock factor of $\cH_h$.
Once the heights include the fixed word's light cone, adding empty
descendants identifies these updates under the compatible embeddings.
We henceforth suppress the height indices, writing $v_P=u^*Pu\vac$;
a direction inside $q$ always comes from the child level.

\begin{lemma}[Ordinary root updates]\label{lem:ordinary-updates}
Ordinary finite words have smooth root limits and smooth adjoint
limits at every sufficient finite height.  Starting at $u=1$,
the limiting left updates for a mixer of angle $\beta$ and a cost of angle $t$
are, respectively,
\begin{equation}\label{eq:ordinary-root-updates}
 u\longmapsto e^{-\ii\beta X}u,
 \qquad
 u\longmapsto e^{-\ii tZq(v_Z)}u.
\end{equation}
The cost update leaves $v_Z$ unchanged.  For a cost conjugated
by a global mixer, replace $Z$ by the same rotated Pauli $P$
both at the root and in $v_P$.  A Gaussian seed gate has the
additional update $u\mapsto e^{-\ii tZs}u$.

All these assertions are unchanged when the top root has $D+1$
children and the descendants have $D$ children, still using
the coefficient $D^{-1/2}$. For an ordinary word $W$, its limiting
insertion $v_Z$ therefore satisfies $v_W=\cE(v_Z)$, with
$\cE$ defined in Proposition~\ref{prop:energy}.
\end{lemma}

\begin{proof}
At finite $D$, factoring out the new child costs gives the exact
left update
\[
 \exp\left(-\ii tZ D^{-1/2}\sum_{j=1}^D
                            U_j^*Z_jU_j\right).
\]
The one-copy observable is centered, because $Z$ is odd under
the joint spin/seed symmetry and the preceding word preserves
that symmetry.  Its vacuum column is $v_Z$.  On each root
$Z$-eigenspace, the update is a product of small one-child
unitaries.  Lemma~\ref{lem:unitary-clt} gives precisely the
cost in \eqref{eq:ordinary-root-updates}; the scalar correction
is zero for an exponential of a fixed observable.  The smooth
estimates follow inductively from the same lemma and
Lemma~\ref{lem:weyl-bounds}.  Mixers are local, and a seed
gate is multiplication by a two-by-two unitary; its input
derivatives have polynomial seed bounds.  Both therefore
preserve the smooth domain.  The update commutes with the
root $Z$, proving the assertion about $v_Z$; conjugation proves
the rotated-axis statement.

For the full-root assertion use $M=D+1$ in the occupation
decomposition while retaining $a/\sqrt D$ in each one-copy
factor.  In the proof of Lemma~\ref{lem:unitary-clt}, every
limiting coefficient is multiplied by $M/D\to1$, and every
bound is unchanged up to a constant.  Induction on the fixed
word proves that full-root and cavity-root maps have the same
limit. Padding the subtree by the finite light-cone radius
justifies all factorizations before limits are taken.

The finite-degree ordinary insertion vectors meet the hypotheses of
Proposition~\ref{prop:energy}: parity makes them centered, the
construction preserves rooted symmetry, and the smooth weight
dominates $1+N$. Their edge pairing is the expectation of $Z_iZ_j$.
The root-shift formula and Proposition~\ref{prop:tree-transfer}
therefore give $v_W=\cE(v_Z)$.
\end{proof}

\subsection{Substitution of commuting controls}
\label{sec:commuting-controls}

The ordinary cost update uses a scalar time multiplying one field.
A controlled operation instead inserts a commuting list of signal
operators into the parameters of a local matrix. The following lemma
makes that substitution legitimate,
including its derivatives and adjoints; no choice of independent
Gaussian coordinates is needed.

\begin{lemma}[Smooth commuting substitution]
\label{lem:commuting-substitution}
Let $S_D=(S_{D,1},\ldots,S_{D,b})$ be a fixed finite list of
strongly commuting self-adjoint operators.  Suppose their
polynomials converge smoothly and their joint exponentials satisfy,
for each $m$, bounds
\begin{equation}\label{eq:joint-exponential-bound}
 \|W_D^m e^{\ii t\cdot S_D}\xi\|
 \leq C_m(1+|t|)^{k_m}\|W_D^{r_m}\xi\|,
\end{equation}
uniformly in $D,t$, and converge strongly on smooth inputs.
Let $F_D(x)$ and $F_D(x)^*$ have uniform smooth bounds and
converge with all their jets on compact $x$-sets, and suppose
$F_D(x)$ commutes with the spectral projections
of every $S_{D,j}$.  Then the substitution $F_D(S_D)$ has uniform
smooth bounds and converges smoothly.  The assertion includes
extra scalar parameters and their jets.  If $F_D(x)$ is unitary
for real $x$, so is $F_D(S_D)$.
The limiting list $S$ consists of the strongly commuting
self-adjoint generators of the limiting joint unitary group.
Their actions and polynomials on the smooth domain agree with
the assumed polynomial limits.

For a real smooth function $f(s,x)$ whose derivatives have
polynomial growth in $x$, uniformly in $s$, the operators
$e^{\ii t f(s,S_D)}$ have polynomial smooth bounds in $t$.
Polynomial growth in a Gaussian $s$ is also permitted with the
seed-moment weight.  In every seed-dependent assertion, the
local seed multiplication algebra is required to strongly
commute with every $S_{D,j}$ and with the operator coefficient
families; thus $f(s,S_D)$ is a joint commuting substitution.
No nondegeneracy assumption on the joint
covariance of the $S_{D,j}$ is needed.
\end{lemma}

\begin{proof}
First identify the limiting signal operators.  Write
$U(t)=\lim_D e^{\ii t\cdot S_D}$ for the strong operator limit.
Strong limits of products of uniformly bounded operators preserve
the group law, and the limits at $t$ and $-t$ preserve inverses;
thus each $U(t)$ is unitary.  For smooth $\xi$,
\[
 \|(e^{\ii t\cdot S_D}-1)\xi\|
       \leq\sum_{\alpha=1}^b|t_\alpha|\|S_{D,\alpha}\xi\|.
\]
The uniform polynomial bounds give strong continuity on this dense
domain, hence on the full space.  The strongly continuous joint
unitary group therefore determines strongly commuting self-adjoint
generators $S_1,\ldots,S_b$.  Passing to the limit in the integrated
one-parameter group identities identifies each generator on smooth
vectors with the corresponding polynomial limit.  The smooth bounds
make this domain invariant, so iteration identifies all polynomial
limits.  This is the realization used in $F(S)$.

We give the estimates, which also specify the substitution on the
smooth domain when $F_D$ is unbounded. For such coefficients, commutation
with the spectral measures is understood in the closed-operator
(or decomposable-operator) sense. On the invariant smooth domain we use
the equivalent identities with the joint unitary group. Sharp spectral
projections are not required to preserve the smooth domain. These
identities pass to the limit because the moving inputs
$e^{\ii t\cdot S_D}\xi$ converge in every required graph norm, as do
the coefficients acting on them. Thus the limiting coefficients retain
the commutation needed for the Fourier products below. Suppress $D$.  Fix an
output weight $m$, and take $r_m,k_m$ from
\eqref{eq:joint-exponential-bound}.  Choose an integer
$q>k_m+b$.  For a sufficiently large integer $N$, the family
\[
 G(x)=(1+|x|^2)^{-N}F(x)
\]
and its first $q$ derivatives are integrable in $x$ as operators
from some fixed input graph norm to the $r_m$th output graph
norm.  Such an $N$ exists because only finitely many polynomial
growth exponents are involved.  With
$\widehat G(t)=(2\pi)^{-b}\int e^{-\ii t\cdot x}G(x)\,\dd x$,
integration by parts gives
\[
 \|W^{r_m}\widehat G(t)\eta\|
       \leq C(1+|t|)^{-q}\|W^r\eta\|.
\]
Therefore
\begin{equation}\label{eq:operator-fourier-substitution}
 F(S)\xi=\int_{\R^b} e^{\ii t\cdot S}\widehat G(t)
                (1+|S|^2)^N\xi\,\dd t
\end{equation}
converges absolutely in the $m$th graph norm and has a bound of
the form \eqref{eq:smooth-family-bound}.  The polynomial on the
right acts first; its bound can be absorbed by increasing the
input weight last.  This order of choices avoids any circular
dependence of exponents.

Integration by parts in $t$, with the same integrable graph-norm
bounds, shows that the result does not depend on $N$: multiplying
a symbol by $1+|x|^2$ corresponds to applying $1-\Delta_t$ to
its Fourier transform and hence to multiplication by
$1+|S|^2$ after integration.  The common smooth domain is
preserved by all factors.  Fubini's theorem and the convolution
formula show on that domain that
\[
              F(S)G(S)=(FG)(S).
\]
In this identity the order of $F$ and $G$ is retained; one only
commutes the Fourier factors and polynomials in $S$ past the
operator coefficients, as permitted by the hypothesis.  It can
first be checked for rapidly decreasing smooth symbols and
then for the divided symbols above, moving their polynomial
factors to the input.  Taking adjoints in pairings of smooth
vectors likewise gives $F(S)^*=F^*(S)$ on the common domain.
For scalar symbols this construction is the usual joint
spectral functional calculus by Fourier inversion.  In
particular the constant symbol $1$ substitutes to the identity.
Thus a pointwise unitary symbol and its adjoint substitute to
mutually inverse isometries on a dense invariant domain.  Their
extensions are unitary on the full Hilbert space.  These facts
also identify the construction with the finite-tree substitution
where the controls already have a simultaneous multiplication
representation.

For convergence, first restrict $x$ to a compact set.  Convergence
of the integrands on smooth inputs and the integrable bound just
proved give convergence of the Fourier integrals.  The polynomial
decay supplied by $(1+|x|^2)^{-N}$ controls the discarded tails
uniformly.  The same proof applies to every scalar-parameter
derivative and to adjoints.  Uniform higher output bounds then
give smooth convergence by \eqref{eq:weight-interpolation}.
For coefficients merely measurable in $s$, use dominated
convergence in $s$; for polynomial seed bounds increase the
seed-moment input weight.

Finally, every $x$-derivative of $e^{\ii t f(s,x)}$ is bounded
by a polynomial in $|t|$ times a fixed polynomial in $|x|$
(and $|s|$ in the moment-weight case).  Applying the same
argument proves the last assertion.  The proof used a joint
spectral representation and Fourier inversion, without choosing
independent field coordinates.
\end{proof}

An important supply of lists satisfying the hypotheses is the
following.  If $M_1,\ldots,M_b$ are commuting centered one-copy
signals and the one-copy joint exponentials have polynomial
smooth bounds, apply Lemma~\ref{lem:unitary-clt} to
$K(x)=\exp(\ii\sum x_\alpha M_\alpha)$.  Its scalar correction
vanishes, and hence
\begin{equation}\label{eq:commuting-signal-limit}
 \left(D^{-1/2}\sum_jM_\alpha^{(j)}\right)_{\alpha=1}^b
 \longrightarrow \bigl(q(M_\alpha\vac)\bigr)_{\alpha=1}^b.
\end{equation}
The limiting fields commute: the imaginary part of their Gram
matrix is the vacuum expectation of a commutator of the original
signals, divided by $2\ii$, and is zero.  Their polynomials
converge by Lemma~\ref{lem:centered-sums}.  The last assertion
of Lemma~\ref{lem:commuting-substitution} bootstraps polynomial
exponential bounds for smooth functions of such a list.  Starting
with local seed functions and a single commuting spin axis, this
proves the same statements for every fixed finite feedforward
calculation of commuting signals.

\subsection{Exclusion calculations and controlled gates}
\label{sec:exclusion-calculus}

We now prove the stronger calculus needed for synthesis.  Its
objects are temporary operations, not additional gates in the
final QAOA word.  The purpose of the calculus is to justify
replacing such operations by finite products inside an already
chosen surrounding circuit.

An even \emph{mask} is a measurable subset of $S$ invariant
under $\iota$.  For masks $I,J$ write
\begin{equation}\label{eq:masked-adjacency}
 (A_{IJ}f)_i=\1_I(s_i)D^{-1/2}
       \sum_{j\sim i}\1_J(s_j)f_j.
\end{equation}
Masks and the coefficients attached to them are part of a fixed
finite expression.  They are not optimized after sampling a
tree or its seeds.

\begin{definition}[Exclusion calculation and gate]
\label{def:exclusion-gate}
A finite exclusion calculation uses local seeds, one prescribed
commuting spin axis at each control site, smooth pointwise
functions of finite lists of messages, and normalized sums
\eqref{eq:masked-adjacency}.  An axis may depend on the site's
mask.  A formula is split into finitely many named roles when
different occurrences need different masks.  The following
conditions are imposed.
\begin{enumerate}
\item Every transmitted field is odd under simultaneous seed
reflection and spin flip.  Every pointwise function is smooth
in its aggregate inputs, with polynomial bounds for every
derivative.  Its local-seed coefficients are measurable and
uniformly bounded, unless the seed-moment variant is specified.
\item In every normalized sum, the receiving mask is disjoint
from every role mask occurring below that sum.  Thus all
ancestor receiving types are excluded from the descendant
calculation.  In particular a message sent across a cut does
not use the site on the other side of that cut.
\item A gate chooses a target mask and applies a two-by-two
unitary $V$ at each target, with parameters computed by such
calculations.  The matrix $V$ itself must be smooth in the
aggregate inputs, with polynomial bounds for every matrix
derivative, uniformly in the local seed (or with the stated
polynomial seed bound).  If external parameters are present,
the bounds include all their mixed jets, locally uniformly
in those parameters.  The target mask is disjoint from all control
masks used across an edge in those parameters.  At the target
itself only its seed, not its spin, is used as a parameter.
\item All parameter operators of this one gate belong to the
same commuting control algebra.  The gate, including its
two-by-two matrix, preserves the joint symmetry.
\end{enumerate}
The global gate is the product of these controlled matrices.
Its factors commute: their targets are distinct and none is a
control for another factor of the same gate.  A finite-range
local observable is therefore affected by only finitely many
of them at finite $D$.
\end{definition}

The exclusions make the information across a removed edge depend on
finitely many boundary inputs.  We next show that the resulting
finite recursion respects the weighted limits already proved.  The
replacement assertion is formulated inside an existing circuit,
which is the form required by the synthesis argument.

\begin{theorem}[Smooth gate calculus]\label{thm:gate-calculus}
Fix a finite sequence of exclusion gates, ordinary mixers and
costs, and local seed gates.  Its root interaction-picture maps,
their adjoints, and their root insertion vectors have smooth
limits under the occupation embeddings.  They have uniform
smooth bounds at every sufficient finite height.  The same
limit is obtained at a full root with $D+1$ children.  Hence
their centered insertion energies are given by
Proposition~\ref{prop:energy}.

An exclusion gate may be replaced inside any such fixed
sequence by a sequence of pointwise unitary matrices with the
same parameter calculations, provided every fixed number of
aggregate-input derivatives converges on compact sets and has
uniform polynomial bounds.  Convergence almost everywhere in
the local seed, or convergence in seed measure followed by a
subsequence, is sufficient for bounded measurable local-seed
coefficients.  The resulting root maps, adjoints, and insertion
vectors converge smoothly, and their energies converge.

If local-seed coefficients are uniformly bounded, these bounds
hold conditionally on the root seed, with constants independent
of its value, using $W_0=1$.  Root maps commute with root-seed
multiplication.  Consequently, if a root column $c$ has
conditional $m$th graph norm at most $C_m$, then
\begin{equation}\label{eq:masked-column-bound}
 \|W^m\1_Jc\|\leq C_m\sqrt{\mu(J)}.
\end{equation}
The same statements hold for any further scalar boundary-input
jets of a fixed exclusion calculation.  They also hold,
including adjoint jets, for a finite list of external unscaled
parameters, locally uniformly in those parameters, when the
matrix and signal hypotheses hold jointly for their mixed jets.
\end{theorem}

\begin{proof}
We give both the finite-tree recursion and the analytic
induction.  This is necessary because a global operator-norm
error proportional to the number of vertices would not imply
the assertion.

\paragraph{Boundary inputs.}
Unroll the fixed expression and index its normalized-sum occurrences
by their named roles, retaining repeated occurrences separately.
These indices form the finite boundary-channel list
$\alpha=1,\ldots,b$.  Cut an edge immediately above a branch root.
For an occurrence $A_{I_\alpha J_\alpha}f_\alpha$, replace the
parent's contribution to the normalized sum by a formal real
coordinate $y_\alpha$, retaining the receiving mask $I_\alpha$.
Write $y=(y_1,\ldots,y_b)$; these are substitution variables,
not additional random seeds.  A sender supplies the coordinate
$m_\alpha(y)=\1_{J_\alpha}(s)f_\alpha(y)$, evaluated at that
sender.  The receiving child uses $m_\alpha(y)/\sqrt D$ on its
role mask and ignores the coordinate when that role is inactive.
All outgoing messages that are
requested by the parent are unchanged by the cut: their
formulas exclude the parent's receiving type.  The parameters
of the gate are computed from the finite list of zero-boundary
child signals and these added inputs.  At the branch root this
gives a target matrix $V(y)$, equal to identity on nontargets,
and the outgoing message list $m(y)$.  The same list is
used for each child with its receiving masks imposed there.
Indeed, a component sent to a child excludes that child's
type from its own computation, so it contains no contribution
from that child.  The mask exclusions are needed for exactly
this assertion.

We first derive the factorization in the physical coordinates, before
conjugating by the preceding circuit. Retain the target operations at
distance at most $d$ from the branch root, but compute all their
parameter formulas on the padded tree. Let $G^{[d]}(y)$ be this
truncated operation, and temporarily write $V_{\rm phys}(y)$ and
$m_{\rm phys}(y)$ for the root matrix and outgoing messages just
defined. For $d=0$ the operation is only $V_{\rm phys}(y)$. For
$d\geq1$, partitioning the targets among the root and its children gives
\[
 G^{[d]}(y)=V_{\rm phys}(y)
       \prod_{j\ \mathrm{child}}
       G_j^{[d-1]}(D^{-1/2}m_{\rm phys}(y)).
\]
The same formula list is evaluated in each child with its receiving
masks imposed there. The exclusion rule ensures that the input sent
to a child does not feed back through that child's receiving type.

The zero-boundary child operation is the baseline to factor out. Put
\[
 B^{[d]}=\prod_jG_j^{[d-1]}(0),\qquad
 P=\bigotimes_jU_j,\qquad U=Pu,
\]
where $U$ is the preceding subtree circuit and $U_j$ its child
circuits. For $d=0$ take $B^{[0]}=1$. The residual operation in a
child, expressed in its preceding frame, is
\[
 K_j^{[d-1]}(a)
   =U_j^*G_j^{[d-1]}(0)^*G_j^{[d-1]}(a)U_j.
\]
At the current branch use the same definition,
\[
 K^{[d]}(y)=U^*G^{[d]}(0)^*G^{[d]}(y)U.
\]
Thus $K^{[d]}(0)=1$. This compares two operations on the same
branch; it does not discard the zero-input operation.

Conjugate the physical signals by $P$, writing
$V(y)=P^*V_{\rm phys}(y)P$ and $m(y)=P^*m_{\rm phys}(y)P$.
Equivalently, conjugate each zero-boundary child signal by its
$U_j$, while leaving the root spin in its present frame. Define
$E^{[d]}(y)=P^*(B^{[d]})^*G^{[d]}(y)P$. Factoring the baseline
from the physical identity gives the exact finite-degree recursion
\begin{align}
 E^{[d]}(y)
 &=V(y)\prod_{j\ \mathrm{child}}
            K_j^{[d-1]}(D^{-1/2}m(y)),
                         \label{eq:exclusion-E-recursion}\\
 K^{[d]}(y)
 &=u^*E^{[d]}(0)^*E^{[d]}(y)u.
                         \label{eq:exclusion-K-recursion}
\end{align}
For $d=0$ the first line is $E^{[0]}(y)=V(y)$. For $d=1$ it
multiplies the root matrix by the ratios of the child target
matrices. Notice that $m(0)$ need not vanish: removing the parent
input leaves the root's own signals. Hence $E^{[d]}(0)$ is generally
nontrivial, although every ratio $K^{[d]}$ equals one at zero.

The second line follows by canceling $B^{[d]}(B^{[d]})^*$ between
the two factors. At zero external input the same factorization gives
\[
 G^{[d]}(0)U=B^{[d]}P E^{[d]}(0)u.
\]
The new child circuits are $G_j^{[d-1]}(0)U_j$, so the new root
map is
\begin{equation}\label{eq:exclusion-root-update}
                    u_{\rm new}=E^{[d]}(0)u.
\end{equation}
Only targets within the finite influence distance can respond to a
cut input. Taking $d$ above that distance gives the actual root
update; all more distant zero-input operations are already factored
into the child circuits.

The factors and their parameter arguments commute before
the preceding circuit is applied: targets are disjoint from
controls, and each site supplies only its prescribed axis.
More explicitly, let $\mathcal A_j$ be the abelian von
Neumann algebra generated by all seeds and prescribed
control axes in child branch $j$, excluding target spin
axes.  For every scalar $a$, each truncated gate
$G_j^{[d]}(a)$ commutes with $\mathcal A_j$, since it
changes only target spins and has coefficients in that
algebra.  Consequently
\[
 K_j^{[d]}(a)\in\bigl(U_j^*\mathcal A_jU_j\bigr)'.
\]
Before the final root conjugation by $u$, every component
of $m(y)$ and every zero-boundary parameter sum is
affiliated with the abelian algebra generated by the
root controls and the child algebras
$U_j^*\mathcal A_jU_j$.  Thus $K_j^{[d]}(a)$ commutes
with every spectral projection of every such component,
even when that component uses signals from branch $j$
itself.  The unused components of a common message list
cause no exception.  This proves the exact commutation
required for substituting the operator $m(y)$ into the
scalar family.  All these algebras also commute with
root-seed multiplication, as do the preceding root
maps, so the seed-dependent version of commuting
substitution has its additional commutation hypothesis.
This also proves
that any two descriptions giving the same finite-tree
operation, including its cut versions, have the same root
map.

\paragraph{Zero-boundary signals.}
Assume by induction that the preceding maps $u,u^*$ have
the asserted smooth bounds and limits.  Compute a signal
in its physical control algebra from the bottom roles up.
At a sum, its zero-boundary child columns have already been
computed; after conjugation by the preceding child maps,
they still commute, are still centered by odd symmetry,
and have the required smooth bounds.  A signal at the new
root is calculated pointwise in the resulting child sums
and the root's chosen axis, and is then conjugated by the
preceding $u$ when its one-copy Heisenberg version is needed.

Here is the induction that also controls exponentials.
Local commuting seed/axis variables plainly have joint
exponentials with polynomial smooth bounds.  Suppose this
is true for the child signal list.  Lemma~\ref{lem:unitary-clt}
gives \eqref{eq:commuting-signal-limit}, including the
polynomial frequency bounds on joint exponentials, and
Lemma~\ref{lem:centered-sums} gives polynomial convergence.
Lemma~\ref{lem:commuting-substitution} then computes each
smooth pointwise function and gives polynomial bounds for
the joint exponentials of every new finite signal list.
Conjugating by $u$ and $u^*$ preserves those bounds and
convergence.  This proves the signal assertions in order
of calculation depth, without any assumption on an unknown
outgoing exponential.

\paragraph{Propagation of a gate across a cut.}
Proceed next in $d$ in
\eqref{eq:exclusion-E-recursion}--\eqref{eq:exclusion-K-recursion}.
For $d=0$, $V(y)$ is a smooth two-by-two matrix function of
the signal lists just treated; commuting substitution gives
all assertions.  Suppose they hold for the child family
$K^{[d-1]}(a)$, including all input jets and adjoints.  It
equals identity at $a=0$.  Treating the inputs as odd under
the symmetry gives
$\mathscr P K^{[d-1]}(a)\mathscr P=K^{[d-1]}(-a)$,
where $\mathscr P\vac=\vac$.  Therefore each first derivative
has vacuum mean zero.

For scalar $a$, apply Lemma~\ref{lem:unitary-clt} to the
child product.  With
$\partial_\alpha K^{[d-1]}(0)=\ii T_\alpha$, its limiting
value is
\begin{equation}\label{eq:exclusion-product-limit}
 \exp\left\{\ii\sum_\alpha a_\alpha q(T_\alpha\vac)
 +\frac12\sum_{\alpha,\lambda}a_\alpha a_\lambda
    \left\langle K_{\alpha\lambda}
       +\tfrac12\{T_\alpha,T_\lambda\}\right\rangle_0
 \right\}.
\end{equation}
The second term is scalar for scalar $a$.  After substitution
it can be an operator in the commuting control algebra; it
must therefore be retained.  The product lemma supplies
polynomial bounds for every scalar $a$-jet.  The components
of $m(y)$ are smooth functions of a finite list of commuting
normalized signals and a root control axis.  The finite-tree
commutations already established imply the hypotheses of
Lemma~\ref{lem:commuting-substitution}, which substitutes
$a=m(y)$.  Use the spectral projections of the fixed root
axis if needed.  Multiplication by $V(y)$ completes the
induction for $E^{[d]}(y)$, and conjugation in
\eqref{eq:exclusion-K-recursion} completes it for $K^{[d]}(y)$.
The argument includes all boundary-input derivatives by
the ordinary chain rule.

This proves \eqref{eq:exclusion-root-update} with smooth
bounds and limits.  It adds one gate to the preceding
sequence.  Ordinary cost and mixer updates may be inserted
between these gates by Lemma~\ref{lem:ordinary-updates};
its proof requires only smooth bounds on the preceding maps
and preservation of symmetry.  Local seed gates are handled
as in that lemma.  Induction on the finite sequence proves
the first assertion.  The same proof with $D+1$ top factors
instead of $D$ gives the same limit, since $(D+1)/D\to1$.
Smoothness supplies the bounded occupation required for
Proposition~\ref{prop:energy}.

External parameters can be carried through this induction
without being scaled.  At the child-product step use the
unscaled-parameter assertion of Lemma~\ref{lem:unitary-clt};
at the substitution step use the mixed-jet version of
Lemma~\ref{lem:commuting-substitution}.  Each unitary strong
limit is unitary, so adjoint convergence follows from the
identity in the proof of Lemma~\ref{lem:unitary-clt}; its
jet version follows by finite differences with the adjoint
jet bounds.  For real commuting signals the adjoint is the
same signal, and compositions use the reversed products of
the already convergent adjoint families.  Thus adjoint
convergence is established at each induction step, rather
than inferred from bounds alone.

\paragraph{Continuity inside a sequence.}
Repeat these two inductions with a further approximation
index.  Compact convergence of a fixed number of matrix
derivatives and uniform polynomial bounds are precisely
the hypotheses of commuting substitution.  The normalized
sum and product lemmas allow the child maps themselves to
vary with this index.  Hence every new root map and its
adjoint converge on smooth vectors with uniform smooth
bounds.  A fixed finite surrounding sequence is handled
by induction and the product rule.  The assertion does
not make a bound uniform in the length of that sequence.
For a pointwise product formula on a single target and a
common commuting control algebra, the entire product is
first regarded as one matrix-valued function; the theorem
then applies to its smooth matrix limit.  This is how
arbitrarily long local commutator or sum formulas can be
used without an estimate proportional to graph size.
For measurable seed limits, dominated convergence in the
root seed and then the same child induction give the
stated extension.  Norm continuity
\eqref{eq:energy-continuity} proves energy convergence.

\paragraph{Conditional bounds and masks.}
Take $W_0=1$.  Every root map commutes with root-seed
multiplication, by inspection of the finite-tree gates
and their factorizations.  Disintegrate the root Hilbert
space over that seed.  All preceding estimates remain
valid in each fiber with the same constants: coefficients
are uniformly bounded there, and the normalized child
columns are averaged over independent child seeds.
Moreover the weight commutes with root-seed multiplication.
Integrating the squared conditional estimate over $J$
proves \eqref{eq:masked-column-bound}.  With polynomial
seed coefficients use instead the seed-moment weight and
integrate the polynomial majorants.  This completes the
proof.
\end{proof}

The gate theorem controls a whole unitary.  In the helper
contractions we also need its first variation with respect to time:
this is a Hamiltonian response to boundary inputs.  The next
corollary derives that response from the same unitary limit, so that
the scalar second-order term and lower-level responses are retained.

\begin{corollary}[Additive cut responses]\label{cor:additive-response}
Fix a preceding finite sequence as in Theorem~\ref{thm:gate-calculus}.
Suppose a temporary Hamiltonian is diagonal in a common
collection of site axes, with coefficients given by an
exclusion calculation.  Require its local Hamiltonian
matrices and every aggregate-input derivative to have
polynomial bounds as in Definition~\ref{def:exclusion-gate}.
In a cut branch let $H_D(a)$ denote
the sum of the Hamiltonian terms affected by scalar
boundary input $a$, truncated beyond its finite influence
distance and conjugated by the preceding subtree circuit.
Assume explicitly that
\begin{equation}\label{eq:additive-response-parity}
                  H_D(-a)=\mathscr P H_D(a)\mathscr P,
                  \qquad \mathscr P\vac=\vac,
\end{equation}
where $\mathscr P$ is the joint spin/seed symmetry.
The first two boundary derivatives have smooth limits
$H_\alpha=\lim_D\partial_\alpha H_D(0)$ and
$H_{\alpha\lambda}=\lim_D\partial_\alpha\partial_\lambda H_D(0)$,
and $H_\alpha$ is centered.  These limits concern the response
derivatives; no limit for the extensive $H_D(0)$ is required.
The change in its $D$ child branches when zero boundary input is
replaced by $a/\sqrt D$ is
\[
 \sum_{j=1}^D
   \bigl[H_D(a/\sqrt D)-H_D(0)\bigr]^{(j)}.
\]
Its smooth limit is
\begin{equation}\label{eq:additive-cut-response}
 \sum_\alpha a_\alpha q(H_\alpha\vac)
 +\frac12\sum_{\alpha,\lambda}a_\alpha a_\lambda
                           \langle H_{\alpha\lambda}\rangle_0.
\end{equation}
The correction is realized as the self-adjoint generator of the
limiting time group of the child cut ratios; the display specifies
its action on the smooth domain.  We use the same realization
after commuting substitution and for the correction generators below.
The same formula applies recursively to lower cut responses
and to their input jets.  It is continuous with the smooth
bounds of Theorem~\ref{thm:gate-calculus}, for a fixed finite
list of channels.  Commuting parameter operators may replace
$a$ under the same exclusions.

In particular, if a sequence of such correction generators
$R_j$ satisfies $\|R_j\xi\|\leq\delta_j\|W^r\xi\|$ with
$\delta_j\to0$, then
\begin{equation}\label{eq:small-generator-error}
 \|(e^{\ii tR_j}-1)\xi\|
       \leq |t|\delta_j\|W^r\xi\|.
\end{equation}
If their exponentials and adjoints also have uniform smooth
bounds, these errors tend to zero in every fixed lower
graph norm and may be inserted in a fixed surrounding
sequence.
\end{corollary}

\begin{proof}
Differentiating \eqref{eq:additive-response-parity} at zero
gives $\mathscr P H_\alpha\mathscr P=-H_\alpha$, proving
centering.  To justify all Hamiltonian-jet bounds without
an estimate on an extensive Hamiltonian, introduce an
unscaled real time $t$ and use the unitary cut ratio
\[
 K_D(a;t)=\exp\bigl(\ii t[H_D(a)-H_D(0)]\bigr).
\]
All current Hamiltonian terms commute before conjugation,
and simultaneous conjugation preserves this property.
In particular $H_D(a)$ and $H_D(0)$ commute.  Thus the
display is exactly the prior-frame ratio of the zero-
boundary gate and the gate with boundary input $a$.
A local matrix $e^{\ii t h(x)}$ and all its mixed $t,x$
derivatives have polynomial bounds from the assumed
bounds on $h$ (use Duhamel's formula for matrix derivatives).
The external-parameter version of
Theorem~\ref{thm:gate-calculus}, which also applies to its
truncated cut recursions, therefore gives uniform smooth
bounds and convergence of $K_D(a;t)$ with all time and
boundary-input jets.  In particular
\[
 H_D(a)-H_D(0)=\frac1\ii\partial_tK_D(a;0),\qquad
 \partial_a^\alpha H_D(a)
   =\frac1\ii\partial_a^\alpha\partial_tK_D(a;0)
       \quad(|\alpha|\geq1).
\]
These are the required bounds for the response jets;
no bound for the possibly extensive $H_D(0)$ is asserted
or used.

For the child-product limit compute the first two boundary
jets at zero:
\[
 K_\alpha(t)=\ii t H_\alpha,\qquad
 K_{\alpha\lambda}(t)=\ii t H_{\alpha\lambda}
                -\frac{t^2}{2}\{H_\alpha,H_\lambda\}.
\]
In Lemma~\ref{lem:unitary-clt} one has $T_\alpha=tH_\alpha$;
the two anticommutator terms cancel exactly.  The product
of the $D$ child ratios therefore tends, with all time
jets, to the exponential of $\ii t$ times
\eqref{eq:additive-cut-response}.  Differentiating at $t=0$
establishes the asserted smooth Hamiltonian limit itself.
The finite-degree child products are unitary groups in $t$.
Their strong limits at positive and negative times retain the
group law and inverses, while convergence of time jets gives
strong continuity.  Their self-adjoint generators consequently
act on smooth vectors by \eqref{eq:additive-cut-response};
no separate choice of closure of the displayed expression is used.
The exclusions and Lemma~\ref{lem:commuting-substitution}
allow operator inputs and preserve the group law and this
generator identification.  Repeating this argument at the
finitely many affected levels proves the recursive assertion.

For $R_j$ with this self-adjoint realization, the spectral inequality
$|e^{\ii tx}-1|\leq|t x|$ proves
\eqref{eq:small-generator-error}.  Uniform smooth
bounds and \eqref{eq:weight-interpolation} then give
higher-norm convergence; applying the same reasoning
to the adjoints and using finite composition proves
the final statement.
\end{proof}

\begin{remark}[Scope of uniformity]\label{rem:calculus-uniformity}
The constants in the calculus are uniform in degree and
in the stated pointwise approximation limits, for a
fixed finite calculation and fixed surrounding sequence.
They are not automatically uniform when the number of
roles or channels tends to infinity.  In a refinement
argument, such as the helper contraction below, that
additional uniformity must be proved from estimates on
the combined columns or kernels.  Once it is proved,
the same smooth-domain and continuity arguments apply.
Scalar phases depending only on the root seed commute
with root observables and do not change insertion
vectors or the child conjugations in the cut recursion.
\end{remark}
